\subsection{Joining the one terminal union}

We now join the regular departure to the terminal estimates. Each
terminal union includes its constant-sign connections, so the result
counts $K$ on the whole physical interval.

\begin{proposition}[The logarithmic-ratio count on the physical interval]
\label{rw:anchored-unions}
Assume \eqref{rw:counting:forcing}, and let
$J=[p_0,p_\gamma]$, with
$p_0\in\{p_\sharp/8,p_\sharp/4\}$, be the physical interval
of Proposition~\ref{rw:saturation}. Then
\begin{equation}\label{rw:merge:two}
 Z(K;J)\le2.
\end{equation}
If \eqref{rw:merge:strong-entry} holds at $p_0$, then
\begin{equation}\label{rw:merge:one}
 Z(K;J)\le1.
\end{equation}
All zeros are counted with their full analytic multiplicities.
These assertions verify hypothesis \textup{(III)} of
Section~\ref{rw:counting}.
\end{proposition}
\begin{proof}
Write $J=R\cup T$ by Proposition~\ref{rw:ordered-cover}, with at
most one terminal union and every internal constant-sign connection
retained. Section~11 gives its domains and fixed-fibre forcing.
Use the seam and center convention of Section~\ref{rw:anchors:primary}.
If there is no terminal part, Proposition~\ref{rw:regular-chain}
already gives both conclusions.
For each direct derivative row below,
Lemma~\ref{rw:anchors:calculus} gives two tail zeros, or one after
a regular zero or the supplied strong entrance: Proposition
~\ref{rw:regular-chain} gives $KK_p>0$ at that seam.
A direct bound $Z(K;T)\le1$ gives the same union conclusions;
a constant-sign piece contributes no zero.

The following table records the terminal estimates used in the proof.
Every derivative numerator has a positive multiplier on its actual
passing interval.
\begin{center}\small
\begin{tabularx}{\linewidth}{@{}p{.42\linewidth}X@{}}
\toprule
 Terminal description & Estimate or entrance result used\\ \midrule
 Two finite lower roots
       &$Z(K')\le1$; \ref{rw:finite:two-root-count}\\
 One incoming loss at a finite saddle
       &$K<0$ if $r\ge1$, or $Z(K')\le1$ if $r<1$;
                         \ref{rw:finite:upper-loss}\\
 Strict or repeated one-infinite lower contact
       &$Z(K')\le1$; \ref{rw:infinity:strict-lower-loss},
                         \ref{rw:infinity:repeated-lower-loss}\\
 Signed horizontal finite root or cubic end
       &$Z(K')\le1$; signed-horizontal calculation in
          Section~\ref{rw:infinity:signed-horizontal}, or
          Proposition~\ref{rw:infinity:cubic}\\
 Strict focus or nonjoint opposing tail
       &$Z(K)\le1$; \ref{rw:infinity:strict-focus},
                         \ref{rw:infinity:opposing}\\
 Joint contact and separated tail
       &$Z(K')\le1$ on the whole union;
                         \ref{rw:joint:one-critical}\\
 Finite saddle with regular lower legs
       &Anchored result \ref{rw:anchors:regular-saddle}\\
 Primary three-saddle
       &Fixed-seam union result \ref{rw:anchors:three-saddle}\\
 Regular outgoing lower arc at a primary weak upper end
       &$Z(K')\le1$; \ref{rw:infinity:weak-regular}\\
 Primary weak contact or weak center
       &Union result \ref{rw:anchors:primary-weak}\\
 Secondary finite saddle or weak end
       &\ref{rw:secondary:saddle-anchor},
                         \ref{rw:secondary:weak-anchor}\\
 Other secondary end
       &$Z(K)\le1$ or a strict derivative;
                         \ref{rw:secondary:regular},
                         \ref{rw:secondary:real}\\
\bottomrule
\end{tabularx}
\end{center}
The constant-sign lower or upper slow pieces are those of
Lemma~\ref{rw:finite:incoming-contraction} and
Propositions~\ref{rw:infinity:negative-regular},
\ref{rw:infinity:slow-domination}, together with their
reflections and secondary lower-double overlays.
The joint estimate applies to its whole contact--separated union.
At a cubic end the multiplier is $F_+(t)F_-(f(t))>0$ on passage.
It remains to treat interior connections and the entrance-dependent
anchors.

\emph{A lower loss before a finite saddle.}
Suppose an incoming double root precedes a finite upper
saddle. Write its lower parameters as
$\bar B=b_0^2,\bar a=2b_0$ and the upper eigenvalues
as $\alpha_0,-d_0$, with $\alpha_0,d_0>0$.
The root inequality $1/b_0<q_c=1/d_0$ gives
\[
 \tau_1=2b_0+\alpha_0-d_0>b_0+\alpha_0>0 .
\]
The regular action difference is positive at small input
and tends to $+\infty$ at the incoming lower root.
The regular one-simple-zero result therefore makes
$K<0$ throughout the intervening limiting regular
interval. Compact convergence gives its actual sign
on the retained overlaps, and incoming contraction
continues it through the finite connection.
There was no regular zero. The remaining saddle interval
has at most two zeros of $K$ by
Proposition~\ref{rw:finite:upper-loss}; for $r\ge1$
it has none. Hence the complete union has at most two.
If the strong entrance held, $\tau_1>0$ would force
$\sigma_1<0$ and a positive initial value of $K$,
contradicting the negative regular sign.
Thus this configuration cannot occur under that entrance.
The coincident incoming loss needs no intervening
connection and its derivative bound is applied directly.
Reflection proves the opposite orientation with its
corresponding interval and signs.

\emph{A lower root before or after a one-infinite threshold.}
A simple root before the threshold ends passage and
cannot be followed by a second counted part.
For a double outgoing root $r$ strictly before a strict
or repeated upper threshold $P$, the inverse incoming
contraction estimate gives $K>0$ through its finite
connection. The outgoing action diverges while the
incoming action there is regular; the small-input sign
is positive as well, so the regular one-zero result
excludes a preceding regular zero.
For example at the repeated field
$\bar a=-2/r,\bar B=r^{-2},\bar c=2/P,
\bar m=-P^{-2}$, $r<P$ gives
$\tau_1<0$ and $\sigma_1>0$. At a strict threshold
write the positive roots of the upper characteristic
polynomial as $\gamma_+>\gamma_->0$.
The inequality $r<1/\gamma_+$ gives
$\tau_1=\gamma_++\gamma_--2/r<0$ and
$\sigma_1>0$ again.
Thus the strong entrance would require a negative
initial value and is excluded.
The upper contact and separated tail after that connection
have the same strictly negative $K_p$ on their open
overlap: this is the double-before-threshold numerator
in Section~8 and, at a repeated tail, the sign in
\eqref{rw:infinity:opposing-slope}.
Their union has at most one zero. Its preceding positive
connection does not increase this bound.

At a repeated upper threshold with a first lower root
$r>P$, the intervening returning tail is the
constant-sign piece $K<0$ of
Proposition~\ref{rw:infinity:negative-regular}.
The finite-root derivative and its infinite limiting
description are alternative overlapping charts of the
same end, as in Proposition~\ref{rw:ordered-cover}.
For a first simple root
$\bar a=-1/r-\bar Br$ and $\bar Br^2\le1$; hence
\[
 \sigma_1=2P^{-1}(\bar B-P^{-2})<0,\qquad
 \tau_1=2/P-1/r-\bar Br>0.
\]
The same inequalities hold at a double root $r>P$.
The small-input value and the last regular value are
both negative, so again the regular prefix is zero-free.
The strong entrance would have positive sign and is
excluded. The final finite-root piece has a strict
derivative, or the signed horizontal/cubic numerator
has at most one zero. It consequently contributes at
most two zeros, with no count from the preceding
constant-sign connection. Root coincidence with $P$
is one contact case already listed above.

The weak boundary with a lower double strictly before
its threshold uses its constant-sign connection and the
weak-regular numerator of
Proposition~\ref{rw:infinity:weak-regular}.
For an outgoing double, $r<P=1/\bar c$ gives
$\tau_1=\bar c-2/r<0$, $\sigma_1=\bar c/r^2>0$
and a positive initial/connection sign.
For an incoming double, $\bar a=2/r$ gives
$\tau_1=\bar c+2/r>0$ and $\sigma_1=\bar c/r^2>0$;
the incoming action gives the negative connection sign,
and the strong coefficient condition is itself impossible.
In either case the primary
regular prefix has no zero, its value is incompatible
with the strong entrance, and the final numerator has
at most one zero. This gives at most two zeros on the
union. The weak contact exactly at the threshold is
the anchored weak case below.

\emph{Primary anchors.}
Apply the three primary entrance propositions listed in the table.
Coefficient transport supplies their fixed-fibre forcing, including
the primary order-one $\nu$ term; the three-saddle proposition derives
its stronger \eqref{rw:finite:three-forcing} from the same outer cycle.
Their union counts already include the regular prefix and exclude
the incompatible entrances in the two-critical three-saddle and
positive-$\bar B$ weak branches.

\emph{Secondary anchors.}
The complete regular chain supplies
\eqref{rw:secondary:anchor-entry}, either after its zero or from
\eqref{rw:merge:strong-entry}. Together with coefficient transport,
these are the saddle-anchor hypotheses. At the weak end, reflection
to $C>0$ gives $K,K_p<0$, the weak-anchor entrance. The two
propositions give the asserted union counts, retaining lower-double
overlays and their incompatible signs. The other secondary rows use
the table; their regular/focus pieces have already been counted in $R$.

Exhaustion in Section~11 completes the proof: a regular zero leaves
at most one terminal zero; without one the union has at most two.
Under \eqref{rw:merge:strong-entry} the chain is zero-free, each
compatible terminal case leaves at most one zero, and the sign
arguments above exclude the remaining connections.

Physical truncation restricts the counted interval. The endpoint
signs used above belong to the stated comparison continuations.
Under reflection the original central logarithm transforms as
$K_{\rm ref}(q)=-K(p)$ and
$K_{\rm ref}\partial_qK_{\rm ref}=KK_p/q'(p)$.
Thus multiplicities and outward departure are preserved at the
corresponding seam and threshold, while both generator signs reverse.
The reflected outgoing label is the original incoming label.
All noncompact counts use the finite-selection convention of Section~3.
\end{proof}
